\chapter{央行不是印钞厂，是水闸和温度计}
\chaptermeta{08}

\begin{ledetext}
上一章证明了钱主要是商业银行造的。那新闻里天天出现的央行到底在管什么？它管的不是水量，是水价。
\end{ledetext}

\section{先拆掉那台想象中的印钞机}

说起央行，很多人脑子里是一间轰鸣的车间：崭新的钞票像瀑布一样吐出来，旁边站着个戴白手套的人在按加速键。

这幅画里唯一对的部分是，纸币确实要印。

但它只占全社会货币的百分之四上下（第 7 章那张图：M0 十几万亿元，M2 三百多万亿元）。而且印钞不是"投放货币"，是换包装。陈宁去 ATM 取 800 块，是把银行存款换成现金，社会上的钱一分没多。印钞的节奏是被公众的取现需求牵着走的，不是央行想印多少印多少。

它真正的工作只有一句话。

\begin{pullquote}{}
央行决定不了这个经济体里有多少钱，它决定的是钱的价格。
\end{pullquote}

\section{三重身份，和一份长短不一的任务书}

这个机构有点分裂，因为它同时干着三份工作。

\begin{flowlist}
  \flowitem{01}{发行的银行}{它垄断法定货币的发行权。你钱包里那张纸的信用，来自这个机构和它背后国家的征税能力。这是最古老、现在也最不吃力的一份工作。}
  \flowitem{02}{银行的银行}{所有商业银行都在央行开户，账上那笔钱就是第 7 章说的准备金。银行之间转账在这本账上清算，银行缺钱来这里借，银行快倒了来这里求救。\textbf{它九成的日常工作在这一层。}}
  \flowitem{03}{政府的银行}{财政部的账户开在央行，税收进来、支出出去都走这里。但有一条红线：多数国家的法律禁止央行直接给财政透支或在一级市场买国债。这条线一旦模糊，货币的锚就松了。中国《人民银行法》也有明确规定。}
\end{flowlist}

它要往哪儿开，取决于被写进法律的那份任务书。各国的任务书长得很不一样，这直接决定了 2026 年它们为什么会走散。

\begin{tablewrap}
\begin{xltabular}{\linewidth}{>{\raggedright\arraybackslash}p{0.20\linewidth}XX}
\toprule
\bthead{央行} & \bthead{法定目标} & \bthead{后果} \\
\midrule
\textbf{美联储} & \textbf{双重使命}：物价稳定加充分就业 & 两个目标打架时（通胀高、失业也高）必须做痛苦取舍，也因此天然容易被政治拉扯 \\
\textbf{欧洲央行} & 以物价稳定为首要目标 & 目标单一、纪律性强，但面对二十个财政不统一的成员国，工具受限 \\
\textbf{中国人民银行} & 保持币值稳定并以此促进经济增长，同时兼顾金融稳定、汇率、就业和结构性任务 & 目标最多，工具箱也最杂。中国这一整套"定向""结构性"工具就是从这儿长出来的 \\
\bottomrule
\end{xltabular}
\end{tablewrap}

目标越多，可用的借口越多，独立性越难守。这不是批评，是结构。

\section{六件工具，六个能推下去的比喻}

\textbf{一、政策利率，总水阀。}这一件最要紧，其余五件都是配角。

美联储的做法叫地板系统。它给银行存在美联储的准备金付息（\term{准备金余额利率}{IORB}\index{090@准备金余额利率}），又给货币基金这类非银机构开了个隔夜逆回购窗口（\term{隔夜逆回购}{ON RRP}\index{020@隔夜逆回购}）。前者大致构成利率上沿，后者构成下沿，市场利率被夹在中间，这条通道叫\term{利率走廊}{interest rate corridor}\index{032@利率走廊}。

\begin{egbox}{举个栗子：两个报价怎么关住一整个市场}
2026 年 8 月，联邦基金利率目标区间是 3.50\% 到 3.75\%。ON RRP 利率在区间下沿附近，IORB 在区间内靠近上沿。

\textbf{下沿怎么守住。}一家货币市场基金账上趴着 84 亿美元隔夜资金。如果市场上有人只肯出 3.42\%，它宁可把钱放进美联储的 ON RRP 拿 3.50\%，零风险。所以没人能用明显低于 3.50\% 的价格借到隔夜钱。

\textbf{上沿怎么压住。}一家银行若能在市场上以 3.55\% 借到钱，转手存进美联储拿 IORB（假设 3.65\%），就是 10 个基点的无风险利差。84 亿美元的规模，年化口径 \textbf{840 万美元}。这种套利会立刻把市场利率压回来。

央行不需要投放任何钱。它把两个报价挂在墙上，整个隔夜市场就被夹住了。\hlmark{这就是"管价格"的字面意思。}

\end{egbox}

中国的路子不同，逻辑一样：以 7 天期逆回购利率为主要政策利率，截至 2026 年 8 月是 \textbf{1.4\%}；用中期借贷便利（MLF）投放中期资金；再通过 18 家报价行的 LPR 机制传导到贷款端。还有一件重型工具是存款准备金率，它调的是银行的长期资金成本和流动性，不是第 7 章批倒的那个货币乘数。

\textbf{二、公开市场操作，每天的微调。}央行在市场上买卖债券、做正逆回购，把银行体系的资金量调到让市场利率贴着政策利率跑的位置。这活儿天天干，像给恒温空调做实时补偿。

\textbf{三、QE 和 QT，把手伸到曲线的右边。}政策利率管得住隔夜，管不住 10 年期。于是央行下场大量买入长期国债，压低长端收益率和\term{期限溢价}{term premium}\index{042@期限溢价}，这叫量化宽松；反过来停止再投资、让持仓自然到期缩表，叫量化紧缩。日本央行 2026 年一边加息一边继续削减国债购买，计划到 2027 年一季度降到每月约 2 万亿日元，走的就是后一条路。

\textbf{四、前瞻指引，用嘴巴影响预期。}今天的长期利率，说穿了是市场对未来一连串短期利率的预期的平均。央行只要让市场相信"我未来两年都不加息"，长端利率现在就下来了。一分钱没花，效果已经发生。这是本世纪央行发明的最便宜的工具。

\begin{databox}{2026：沟通范式正在转向}
2026 年，美联储主席凯文·沃什明确表示要\textbf{淡化前瞻指引}。

理由不难猜。当央行反复承诺又反复被现实打脸（比如 2021 年那句"通胀是暂时的"），指引就从降低不确定性变成了制造不确定性，公信力被一次次透支。

沃什同时强调，美联储"不会对任何市场走势背书"，以及"必要和适当时会毫不犹豫地采取行动"。翻译过来：\hlmark{别指望我提前给你剧本，也别指望我托底你的持仓。}

这是过去十几年货币政策沟通方式的一次真转向。对市场的直接后果是，预期的锚变松了，波动会变大。

\end{databox}

\textbf{五、最后贷款人，消防队。}银行的生意天生借短贷长，只要储户同时来取钱，再健康的银行也会倒。所以危机时得有一个不会缺钱的人站出来放款。19 世纪的英国人白芝浩把这件事写成三条铁律，一百五十年没过时。

\begin{chainflow}
\chainnode{只救还有偿付能力的}\chainarrow \chainnode{按惩罚性高利率放款}\chainarrow \chainnode{凭优质抵押品无限量供应}
\end{chainflow}

消防队这个比喻能一路推下去。只救还救得回来的房子，是因为往一堆焦炭上浇水只是浪费水压；收水费，是为了让人别把消防队当免费喷淋用；水管够，是因为火场里最先要扑灭的不是火，是"水不够"这个念头。恐慌只有在"钱管够"这句话被相信的时候才会停。

\textbf{六、宏观审慎和结构性工具，滴灌。}总水阀只有一个，可经济的不同角落干湿不均。于是有了逆周期资本缓冲、房贷首付比例、行业贷款集中度这类\term{宏观审慎}{macroprudential}\index{022@宏观审慎}工具。中国还额外配了一整套定向再贷款：支农支小、科技创新、绿色低碳，给特定方向的贷款提供低成本资金。

这类工具精准，代价是它在干本该财政干的活。\hlmark{决定谁能拿到便宜的钱，是一个分配决策。}

\section{怎么读一张央行的资产负债表}

新闻里天天说"美联储缩表""央行扩表"，听着玄，其实那张表比你家账本还简单。

\begin{bsheet}
  \bsside{QE 之前 · 某央行（亿元）}{%
    \bsrow{持有的国债}{7840}
    \bsrow{对金融机构的贷款}{1620}
    \bsrow{外汇与黄金}{460}
    \bsrowtotal{资产合计}{9920}
  }
  \bsside{负债 + 净值}{%
    \bsrow{银行准备金}{5930}
    \bsrow{流通中现金}{3140}
    \bsrow{政府存款}{610}
    \bsrow{资本}{240}
    \bsrowtotal{合计}{9920}
  }
\end{bsheet}

资产端是它买来的和借出去的，负债端是它欠出去的。QE 就是资产端买入国债、负债端同时长出准备金，两边一起膨胀。QT 反过来，两边一起收缩。

\begin{egbox}{举个栗子：同样买 780 亿国债，结局可以完全不同}
\textbf{情形 A，卖方是一家银行。}央行资产端国债加 780 亿，负债端准备金加 780 亿；这家银行资产端国债减 780 亿、准备金加 780 亿。陈宁和你我的存款：\textbf{一分没动}。只有第一层的钱变多了，广义货币纹丝不动。

\textbf{情形 B，卖方是一家养老基金。}养老基金不能在央行开户，钱得打进它的开户行。于是央行准备金加 780 亿（记在这家银行名下），这家银行资产端准备金加 780 亿、负债端"养老基金的存款"加 780 亿。\textbf{第一层和第二层同时增加}，广义货币真的变多了。

同一句"央行买了 780 亿国债"，两种结局。所以看到"央行大放水"这四个字，先问一句：\hlmark{钱是从谁手里买的，落进了哪一层？}

\end{egbox}

\section{独立性：一个 1970 年代的老问题，2026 年又新鲜了}

为什么全世界要把货币政策交给一群不用选举、任期很长、谁也管不着的技术官僚？

答案叫\term{时间不一致}{time inconsistency}\index{049@时间不一致}。政客的时间尺度是下一次选举，而降息的好处（就业、股市、支持率）来得快，坏处（通胀）来得慢，还常常落到下一任头上。任何把利率交给政治周期的国家，长期都会得到系统性偏高的通胀。1970 年代那笔学费是用两位数通胀和一场深度衰退交的。

把方向盘交给一个不看民调的人，不是因为他更聪明，是因为他敢做那件正确但不受欢迎的事。

2026 年，这个老问题重新变得刺手。美联储主席换成了凯文·沃什，政治层面对货币政策的压力明显上升，而通胀偏偏没回到目标。沃什那两句"不会对任何市场走势背书""必要和适当时会毫不犹豫地采取行动"，一半说给市场听，一半说给白宫听。

\begin{deepbox}{进阶：可以跳过 · 这场争论和作者的倾向}
\textbf{支持独立性的一方：}历史证据很硬。独立性高的央行长期通胀更低、通胀预期更稳，而且没有以更高的失业作为代价。放弃它就是拿几十年的信用去换几个季度的舒服。

\textbf{质疑的一方：}货币政策早就不只是调利率了。QE 会推高资产价格、拉大贫富差距；救谁不救谁是重大的分配选择；结构性工具直接决定哪个行业拿到便宜的钱。这些都是政治决策，却由一群不经选举的人做出，问责在哪里？

\textbf{我的倾向是：}两边说的都对，但它们说的不是同一件事。设定利率这件事应当独立，证据足够充分，不该动。而 QE、定向再贷款、救助名单这些带明确分配后果的操作，应当由财政和立法机构明确授权、事后被审查。把这两类权力混在一起打包成"央行独立性"来辩护，才是过去十五年真正的漏洞。

\end{deepbox}

\section{2026 年 8 月：四家央行走散了}

过去二十年我们习惯了一件事：美联储一转向，全世界跟着转。2026 年这个默认设定失效了。

\begin{tablewrap}
\begin{xltabular}{\linewidth}{>{\raggedright\arraybackslash}p{0.20\linewidth}XXr}
\toprule
\bthead{央行} & \bthead{2026 年 8 月的位置} & \bthead{方向} & \bthead{关键利率} \\
\midrule
\textbf{美联储} & 2026 年已连续 5 次按兵不动（截至 7 月），FOMC 称通胀"仍然高企"，市场甚至讨论过加息 & 观望偏鹰 & 3.50\%–3.75\% \\
\textbf{欧洲央行} & 2026 年 6 月 11 日上调存款便利利率，三年来首次加息 & 转向收紧 & 已上调 \\
\textbf{日本央行} & 2026 年 6 月 16 日以 7 比 1 加息 25 个基点，为 1995 年 9 月以来最高，同时继续削减国债购买 & 正常化推进中 & 1.0\% \\
\textbf{中国人民银行} & 基调"适度宽松"，强调"精准有效""适时调整"，支持科技创新、提振消费、民营小微、绿色转型 & 宽松但克制 & 1.4\% \\
\bottomrule
\end{xltabular}
\end{tablewrap}

\begin{statrow}{3}
  \statitem{1.0}{\%}{日本央行政策利率，1995 年 9 月以来最高}
  \statitem{1.4}{\%}{中国 7 天期逆回购政策利率}
  \statitem{3.5}{\%}{中国 5 年期以上 LPR（2026 年 8 月 20 日）}
\end{statrow}

一个在收紧，一个在正常化，一个在观望，一个在宽松。

利差就是资本的坡度。日本的利率从负 0.1\% 一路爬到 1.0\%，那些年借便宜日元去买全球资产的套利盘就得重新算账；欧洲开始加息而美国不动，欧美之间的资金流向和汇率就得重新找平衡点。

\begin{pullquote}{}
央行不同步，本身就是资本流动和汇率波动的源头。世界上不再有一个统一的"钱的价格"了。
\end{pullquote}

\section{两个必须扔掉的直觉}

\begin{mythbox}{常见误解}
\mvfrow{误解}{"央行想让经济好，就能让经济好。降息、放水，经济自然起来。"}
\mvfrow{事实}{央行只能改变钱的价格和可得性，它\textbf{替不了任何人做借钱的决定}。传导链有四个地方会断：利率降到接近零仍无人肯借（流动性陷阱）；企业和家庭忙着还债，目标从赚钱变成减债（资产负债表衰退）；银行资本不足或风险偏好塌了，宁可买国债也不放贷；以及最朴素的一种，根本没有值得投的项目。\hlmark{你能把马牵到河边，逼不了它喝水。}}
\end{mythbox}

\begin{mythbox}{常见误解}
\mvfrow{误解}{"降息是利好，说明经济要转好了。"}
\mvfrow{事实}{这个说法是错的，而且通常刚好反过来。\textbf{央行降息，是因为它看见了需要被救的东西。}降息是对坏消息的反应，不是好消息本身。历史上多次深度衰退的起点，都伴着一轮急促的降息。该盯的从来不是"降了没有"，是"为什么降、降得多急"。连续大幅降息往往意味着某个地方在着火。}
\end{mythbox}

\begin{warnbox}{小心}
把央行动作当买卖信号，是散户最常见的亏钱方式之一。政策利率是\textbf{已经发生}的事，资产价格交易的是\textbf{预期}。降息公布那一刻，市场可能早把三次降息算进价格里了，靴子落地反而是利空。本书不提供任何交易建议，只提醒你：新闻标题和你的账户之间，隔着一整条下一章要拆的传导链。

\end{warnbox}

\bookbreak

\begin{takeaway}{本章一句话}
央行管的不是钱的数量，是钱的价格。它是水闸和温度计，不是印钞厂。

\begin{itemize}
  \item 三重身份：发行的银行、银行的银行、政府的银行。九成的日常工作在第二层。
  \item 工具箱六件：政策利率、公开市场操作、QE 与 QT、前瞻指引、最后贷款人、宏观审慎与结构性工具。
  \item 2026 年四家央行走散：欧洲加息、日本升到 1.0\%、美联储停在 3.50\%–3.75\%、中国维持 1.4\% 的适度宽松。不同步本身就是波动的来源。
  \item 降息通常是坏消息的反应。央行能把马牵到河边，逼不了它喝水。
\end{itemize}

\end{takeaway}

\begin{bridgebox}
\textbf{下一章}央行拧的是最上游那个阀门。这股水要经过几道弯才能流到陈宁每月 15 号被划走的那笔钱上？我们把管道一节一节拆开。
\end{bridgebox}

